\documentclass[10pt,a4paper]{amsart}
\usepackage[margin=19mm]{geometry}
\usepackage{amsmath,amssymb,mathtools}
\usepackage[scr=rsfs,cal=euler]{mathalfa}
\usepackage{xfrac}
\usepackage[unicode,hidelinks,bookmarks=false]{hyperref}
\newcommand{\ie}{i.e.\ }
\newcommand{\cf}{cf.\ }
\newcommand{\resp}{respectively\ }
\newcommand{\Subsection}{\subsection}
\providecommand{\nobreakdash}{\nobreak\hbox{-}}
\setlength{\emergencystretch}{2em}
\multlinegap0pt
\usepackage{float}
\usepackage{mathtools}
\usepackage{xcolor}
\usepackage{tikz}
\usepackage{algorithm}\usepackage{algpseudocode}
\usepackage{amssymb}
\usepackage{breqn}
\usepackage[shortlabels]{enumitem}
\newtheorem{claim}{Claim}
\title{On the maximum product of distances of diameter $2$ point sets}
\author{Stijn Cambie, Arne Decadt, Yanni Dong, Tao Hu, Quanyu Tang}
\date{Source: arXiv:2603.07088v1 (CC BY 4.0)}
\begin{document}
\maketitle

\noindent\emph{Source and licence.} On the maximum product of distances of diameter $2$ point sets. Version arXiv:2603.07088v1 (CC BY 4.0), distributed by arXiv under the Creative Commons Attribution 4.0 International licence, http://creativecommons.org/licenses/by/4.0/, as recorded in the arXiv licence metadata for this exact version and shown on the abstract page https://arxiv.org/abs/2603.07088v1. Full licence notice: \texttt{licenses/arxiv-2603.07088-CC-BY-4.0.txt}.\par\smallskip

\noindent\textit{Editorial scope.} Counted unit: the article's claim claim:solve\_equations\_vv1 (source0.tex lines 2484--2492). The article's complete original proof of it is delivered with it (source0.tex lines 2494--2588, one \texttt{proof} environment of 95 lines). No mathematics is written, repaired or re-derived: the statement and the proof are the article's own lines, in the article's own notation, and no retained line is substituted or re-flowed.  Retained uncounted as the source-local context the proof needs: lines 2477--2482.  Nothing was omitted from any retained span, so the package carries no omission record.  No retained line contains a citation, so the wrapper carries no bibliography and no bibliography entry.  No retained line contains a figure environment, an image-inclusion command or drawing code, so no image is delivered and the package declares no asset dependency.  Editorial formatting only, in addition: an AMS \texttt{amsart} preamble that restates the article's own declarations and macros used by the retained lines, namely the environments \texttt{claim} on separate counters, together with the standard packages those lines need.  The retained environments print in this document as Claim 1 in order of appearance, whereas the article prints the same items with the numbering of its own full document.  The complete native source of the article as distributed in the arXiv e-print for this exact version is bundled unchanged as \texttt{sources/195917/source0.tex}.
\begin{align}
\frac1{z_1-2}+\frac1{-2}+\frac1{z_3-2} &= -2\lambda_{2,4}, \label{eq:P4-k4}\\
\frac1{z_1}+\frac1{z_3}+\frac1{2} &= 2\lambda_{2,3}e^{-i\beta}+2\lambda_{2,4}, \label{eq:P4-k2}\\
\frac1{-z_1}+\frac1{z_3-z_1}+\frac1{2-z_1} &= -2\lambda_{1,3}e^{-i\alpha}, \label{eq:P4-k1}\\
\frac1{z_1-z_3}+\frac1{-z_3}+\frac1{2-z_3} &= 2\lambda_{1,3}e^{-i\alpha}-2\lambda_{2,3}e^{-i\beta}. \label{eq:P4-k3}
\end{align}

\begin{claim}\label{claim:solve_equations_vv1}
Solving \eqref{eq:P4-k4}--\eqref{eq:P4-k3} yields uniquely (up to complex conjugation)
\[
e^{i\beta}=\frac34-\frac{\sqrt7}{4}i,\qquad
e^{i\alpha}=-\frac18+\frac{3\sqrt7}{8}i,\qquad
\lambda_{2,4}=\lambda_{1,3}=\frac34,\qquad
\lambda_{2,3}=0.
\]    
\end{claim}

\begin{proof}[Proof of Claim~\ref{claim:solve_equations_vv1}]
Set $A:=e^{i\alpha}$ and $B:=e^{i\beta}$ so that $z_3=2B$ and $z_1=2(A+B)$. From \eqref{eq:P4-k4} we obtain an explicit expression for \(\lambda_{2,4}\):
\begin{equation}\label{eq:lam24-explicit}
4\lambda_{2,4}
=1-\frac{1}{A+B-1}-\frac{1}{B-1}.
\end{equation}
From \eqref{eq:P4-k1} we similarly get
\begin{equation}\label{eq:lam13-explicit}
4\lambda_{1,3}
=1+\frac{A}{A+B}+\frac{A}{A+B-1}.
\end{equation}
Substituting \eqref{eq:lam24-explicit} into \eqref{eq:P4-k2} and simplifying yields
\begin{equation}\label{eq:lam23-explicit}
4\lambda_{2,3}
=\frac{(A+2B-1)\bigl(A(2B-1)+2B(B-1)\bigr)}{(A+B)(B-1)(A+B-1)}.
\end{equation}


Since the KKT multipliers are real, we must have
\(\lambda_{2,4},\lambda_{2,3}\in\mathbb R\).
Using \(\overline A=A^{-1}\), \(\overline B=B^{-1}\) (because \(|A|=|B|=1\)),
the condition \(\lambda_{2,4}=\overline{\lambda_{2,4}}\) is equivalent (after clearing denominators,
and using \(B\neq 1\) since \(z_3\neq z_4\)) to
\begin{equation}\label{eq:real-lam24-poly}
2A^2B^2 - A^2B - A^2 + 2AB^3 -3AB^2 -3AB +2A -B^3 -B^2 +2B=0.
\end{equation}
Likewise, \(\lambda_{2,3}=\overline{\lambda_{2,3}}\) and \eqref{eq:lam23-explicit} give
\begin{equation}\label{eq:real-lam23-factor}
(A+B^2)\bigl(3A^2B-3A^2+3AB^2-8AB+3A-3B^2+3B\bigr)=0,
\end{equation}
again after clearing denominators.


Assume for contradiction that the second factor in \eqref{eq:real-lam23-factor} vanishes, i.e.
\begin{equation}\label{eq:F-branch}
3A^2B-3A^2+3AB^2-8AB+3A-3B^2+3B=0.
\end{equation}
Rewrite \eqref{eq:F-branch} and \eqref{eq:real-lam24-poly} as quadratic equations in $A$:
\[
3(B-1)A^2+(3B^2-8B+3)A-3B(B-1)=0,
\]
and
\[
(B-1)(2B+1)A^2+(2B^3-3B^2-3B+2)A-B(B-1)(B+2)=0.
\]
Since $B\neq 1$, multiply the first equation by $(2B+1)$ and subtract $3$ times the second equation.
The $A^2$-terms cancel, and a short simplification gives
\[
(B-1)\bigl((4B-3)A+3B(B-1)\bigr)=0,
\]
hence
\begin{equation}\label{eq:A-rational}
A=-\frac{3B(B-1)}{4B-3}.
\end{equation}
Taking absolute values and using $|A|=|B|=1$ yields
\[
|4B-3|=3|B-1|.
\]
Now $|B-1|^2=2-2\cos\beta$ and $|4B-3|^2=25-24\cos\beta$, hence
\[
25-24\cos\beta=9(2-2\cos\beta)\quad\Rightarrow\quad \cos\beta=\frac{7}{6},
\]
a contradiction. Therefore the second factor in \eqref{eq:real-lam23-factor} cannot vanish, and so
\begin{equation}\label{eq:A=-B2}
A=-B^2.
\end{equation}



Substituting \eqref{eq:A=-B2} into \eqref{eq:real-lam24-poly} gives the factorization
\[
0=B(B+1)(B^2-B+1)(2B^2-3B+2).
\]
Now \(B\neq 0\). Also \(B=-1\) is impossible because then
\(|z_3-z_4|=|-2-2|=4>2\).
Finally, \(B^2-B+1=0\) implies \(B=e^{\pm i\pi/3}\), hence
\(|z_3-z_4|=|2B-2|=2\), i.e.\ the edge \(\{3,4\}\) would also be active, contradicting the path assumption.
Therefore
\begin{equation}\label{eq:B-quadratic}
2B^2-3B+2=0.
\end{equation}
Solving \eqref{eq:B-quadratic} gives \(B=\frac34\pm \frac{\sqrt7}{4}i\), and then \eqref{eq:A=-B2} gives
\(A=-B^2=-\frac18\mp \frac{3\sqrt7}{8}i\).
Substituting these values into \eqref{eq:lam24-explicit}--\eqref{eq:lam23-explicit} yields
\[
\lambda_{2,4}=\lambda_{1,3}=\frac34,\qquad \lambda_{2,3}=0,
\]
and the two choices of sign correspond to complex conjugation.



The values of \(A,B,\lambda_{1,3},\lambda_{2,3},\lambda_{2,4}\) obtained above were derived from
\eqref{eq:P4-k4}, \eqref{eq:P4-k2}, \eqref{eq:P4-k1} together with the reality constraints on the multipliers.
Substituting them into \eqref{eq:P4-k3} shows that \eqref{eq:P4-k3} is satisfied as well. This completes the proof.
\end{proof}

\end{document}
